% !TeX root = ../navier-stokes-manuscript.tex
% ============================================================
% 2.1 Vortex Stretching
% ============================================================

The momentum equation has four pieces:
\[
  \underbrace{\partial_t u}_{\text{local acceleration}}
  +
  \underbrace{(u\cdot\nabla)u}_{\text{nonlinear transport}}
  =
  \underbrace{-\nabla p}_{\text{pressure}}
  +
  \underbrace{\nu\Delta u}_{\text{viscous diffusion}}.
\]
The central regularity problem lies in the competition between nonlinear transport and viscosity.
The hope is simple: viscosity smooths the fluid faster than nonlinear transport can sharpen it.
If that remains true at every scale, the velocity stays smooth.

\subsection{Vortex Stretching}\label{sub:ThePerfectlyStretchingVortex}

Three dimensions contain one motion that puts this hope under direct pressure. Take a small material vortex segment and follow the fluid that forms it. Axial strain lengthens the segment while incompressibility contracts its transverse cross-section. The rotating fluid is drawn inward as the core narrows, and the aligned stretching contribution supplies faster rotation. Viscosity acts throughout this motion, so the spin-up describes the nonlinear contribution; net vorticity grows only when that contribution exceeds the simultaneous viscous loss.

Let \(e\) be the unit direction of the vortex axis, let \(L(t)\) be the segment's length, and suppose the strain is locally uniform across it. Writing \(S\) for the symmetric part of the velocity gradient, axial extension at fractional rate \(\sigma(t)\) means
\[
  S:=\frac12\bigl(\nabla u+(\nabla u)^{\mathsf T}\bigr),
  \qquad
  Se=\sigma(t)e,
  \qquad
  \sigma(t)>0,
  \qquad
  \frac{\dot L(t)}{L(t)}
  =
  e\cdot Se
  =
  \sigma(t).
\]
The lengthening fluid cannot keep the same cross-section. Its material volume \(\mathcal V\) is preserved, so define the effective transverse area by \(A(t):=\mathcal V/L(t)\) and the effective radius by \(r_{\mathrm{eff}}(t):=\sqrt{A(t)/\pi}\). Incompressibility then fixes both contraction rates:
\[
  \nabla\cdot u=0,
  \qquad
  \frac{d}{dt}(A L)=0,
  \qquad
  \frac{\dot A}{A}=-\sigma(t),
  \qquad
  \frac{\dot r_{\mathrm{eff}}}{r_{\mathrm{eff}}}
  =
  -\frac{\sigma(t)}{2}.
\]
Thus the radius loses one half of the axial fractional growth rate. This transverse narrowing accompanies the stretching; it does not replace the stretching term that changes the rotation.

Let \(\omega:=\nabla\times u\) be the vorticity. Along the material motion generated by \(u\), it obeys
\[
  \frac{D\omega}{Dt}
  =
  S\omega+\nu\Delta\omega,
  \qquad
  \frac{D}{Dt}
  :=
  \partial_t+u\cdot\nabla,
\]
and when \(\omega=|\omega|e\) is aligned with the stretching axis,
\[
  S\omega
  =
  \sigma(t)\omega,
  \qquad
  \frac12\frac{D|\omega|^2}{Dt}
  =
  \sigma(t)|\omega|^2
  +
  \nu\,\omega\cdot\Delta\omega.
\]
The positive term \(\sigma(t)|\omega|^2\) is the rotational gain supplied by axial stretching as the core contracts. The term \(\nu\,\omega\cdot\Delta\omega\) is present at the same instant and has no prescribed pointwise sign. Rotation strengthens on an interval only when their sum remains positive.

The first available attempt to rule out this narrowing is the finite kinetic energy. The energy identity proved in \Cref{prop:ClassicalKineticEnergyIdentity}, together with the antisymmetric part of \(\nabla u\), gives
\[
  \norm{u(t)}_2
  \leq
  \norm{u_0}_2
  <\infty,
  \qquad
  \left|\frac12\bigl(\nabla u-(\nabla u)^{\mathsf T}\bigr)\right|_{\mathrm F}
  =
  \frac{|\omega|}{\sqrt2}
  \leq
  |\nabla u|_{\mathrm F}.
\]
Finite \(L^2\) velocity does not give a pointwise bound on its gradient. A shrinking region may carry increasing rotation while the total kinetic energy stays finite. The surviving variable is the width of the concentration, and the next test is whether viscosity gains any relative strength when that width is reduced.

\divider{\NavierStokes{} Scaling}

Fix a time \(t_0<T_*\) and write \(\ell:=r_{\mathrm{eff}}(t_0)\) for the effective radius of the material segment at that time. For \(\lambda>1\), compare it with the exact \NavierStokes{} field at width \(\lambda^{-1}\ell\):
\[
  u_\lambda(x,t)
  :=
  \lambda u(\lambda x,\lambda^2t),
  \qquad
  p_\lambda(x,t)
  :=
  \lambda^2p(\lambda x,\lambda^2t).
\]
At time \(\lambda^{-2}t_0\), the corresponding segment in \(u_\lambda\) has effective radius \(\lambda^{-1}\ell\). The field \(u_\lambda\) is another solution viewed at another scale, not a later state of the material segment in \(u\). Its four momentum terms satisfy
\[
\begin{aligned}
  \partial_tu_\lambda(x,t)
  &=
  \lambda^3(\partial_tu)(\lambda x,\lambda^2t),
  \\
  (u_\lambda\cdot\nabla)u_\lambda(x,t)
  &=
  \lambda^3\bigl((u\cdot\nabla)u\bigr)(\lambda x,\lambda^2t),
  \\
  \nabla p_\lambda(x,t)
  &=
  \lambda^3(\nabla p)(\lambda x,\lambda^2t),
  \\
  \nu\Delta u_\lambda(x,t)
  &=
  \lambda^3\nu(\Delta u)(\lambda x,\lambda^2t).
\end{aligned}
\]
Nonlinear transport and viscous diffusion both gain the factor \(\lambda^3\); local acceleration and pressure gain it as well. Finer width gives viscosity no automatic advantage over the sharpening motion. Since the equation preserves their relative strength, the concentration must be measured at a level that also survives this change of scale.

\divider{Critical Height}

The familiar norms do not survive it together. Let \(\Lambda:=(-\Delta)^{1/2}\), so that \(\norm{u}_{\dot H^s}=\norm{\Lambda^s u}_2\). The Fourier transform of the rescaled velocity is
\[
  \widehat{u_\lambda}(\xi,t)
  =
  \lambda^{-2}\widehat u(\lambda^{-1}\xi,\lambda^2t),
\]
and changing variables in the frequency integral yields
\[
  \norm{u_\lambda(t)}_{\dot H^s}^2
  =
  \lambda^{2s-1}
  \norm{u(\lambda^2t)}_{\dot H^s}^2.
\]
At \(s=0\), kinetic energy loses one power of \(\lambda\). At \(s=1\), the first-derivative norm gains one. The exponent \(2s-1\) selects \(s=\tfrac12\) as the level between them that does not move. Define the critical height of the velocity field by
\[
  H(t)
  :=
  \frac12\norm{u(t)}_{\dot H^{1/2}}^2
  =
  \frac12\norm{\Lambda^{1/2}u(t)}_2^2.
\]
This is a global functional of \(u(t)\), not the geometric height of the vortex core. Writing \(H_\lambda\) for the functional evaluated on the comparison solution \(u_\lambda\), the three adjacent levels transform as
\[
  \norm{u_\lambda(t)}_2^2
  =
  \lambda^{-1}\norm{u(\lambda^2t)}_2^2,
  \qquad
  H_\lambda(t)=H(\lambda^2t),
  \qquad
  \norm{\nabla u_\lambda(t)}_2^2
  =
  \lambda\norm{\nabla u(\lambda^2t)}_2^2.
\]
Energy can fall and the gradient can rise across comparison scales while \(H\) remains unchanged. What still matters along the original solution is whether nonlinear transport can raise this invariant level faster than viscosity lowers it.

\divider{Nonlinear Production versus Viscous Dissipation}

Pairing the equation with \(\Lambda u\) isolates the signed contribution of nonlinear transport. Define
\[
  \mathcal P_{1/2}[u](t)
  :=
  -\operatorname{Re}
  \pair
    {\Lambda^{1/2}u(t)}
    {\Lambda^{1/2}\bigl((u(t)\cdot\nabla)u(t)\bigr)}_{L^2}.
\]
The pressure pairing vanishes because \(\nabla\cdot u=0\), while the Laplacian contributes \(-\nu\norm{\Lambda^{3/2}u}_2^2\). Hence the critical height evolves by
\[
  H'(t)
  +
  \nu\norm{\Lambda^{3/2}u(t)}_2^2
  =
  \mathcal P_{1/2}[u](t).
\]
The sign is decisive. The height rises only when \(\mathcal P_{1/2}[u](t)\) exceeds the simultaneous viscous dissipation \(\nu\norm{\Lambda^{3/2}u(t)}_2^2\). Under the comparison scaling, the two rates become
\[
  \mathcal P_{1/2}[u_\lambda](t)
  =
  \lambda^2\mathcal P_{1/2}[u](\lambda^2t),
  \qquad
  \nu\norm{\Lambda^{3/2}u_\lambda(t)}_2^2
  =
  \lambda^2\nu\norm{\Lambda^{3/2}u(\lambda^2t)}_2^2.
\]
Both acquire the same square factor, leaving production and dissipation scale-matched. That factor also contracts the comparison time coordinate by \(\lambda^{-2}\), but scale covariance alone does not say how long an actual narrowing event in \(u\) takes. A separate viscous evolution supplies the relevant comparison time.

\divider{The Heat Clock}

Let \(v\) solve the linear heat equation \(v_t=\nu\Delta v\). Over an interval \(\delta t\), each Fourier mode evolves by
\[
  \widehat v(\xi,t+\delta t)
  =
  e^{-\nu|\xi|^2\delta t}\widehat v(\xi,t).
\]
A structure localized at spatial width \(r\) has associated frequency \(r^{-1}\), expressed by \(|\xi|\simeq r^{-1}\). Its viscous change becomes order one when
\[
  \nu|\xi|^2\delta t\simeq1,
\]
so the linear comparison time attached to width \(r\) is
\[
  \tau_\nu(r)
  :=
  \frac{r^2}{\nu}.
\]
This heat time is not a proved lifetime for the material vortex. It says that halving a spatial width quarters the time on which viscosity acts at the associated frequency. More generally,
\[
  \tau_\nu(\lambda^{-1}r)
  =
  \lambda^{-2}\tau_\nu(r).
\]
Successive narrower widths thus carry successively shorter linear comparison times. Their arithmetic can be summed before any claim is made that one nonlinear solution realizes them.

\divider{The Zeno Cascade}

Begin with a reference spatial width \(r_0>0\), halve it repeatedly, and attach to each width its heat comparison time:
\[
  r_n:=2^{-n}r_0,
  \qquad
  \tau_n:=\frac{r_n^2}{\nu},
\]
where \(r_n\) is a concentration width and \(r_n^{-1}\) is its associated frequency. These comparison times satisfy
\[
  \tau_n
  =
  \frac{r_0^2}{\nu}4^{-n},
  \qquad
  \sum_{n=0}^{\infty}\tau_n
  =
  \frac{4r_0^2}{3\nu}
  <
  \infty.
\]
The finite sum is arithmetic, not a realized cascade or a singularity. To connect it to the original velocity field, fix \(0<\theta<1\) and define an actual vorticity concentration width \(\rho(t)\) as the infimum of radii \(r>0\) for which some ball \(B(x,r)\) contains at least the fraction \(\theta\) of the total enstrophy \(\norm{\omega(t)}_2^2\). One \NavierStokes{} solution would then have to pass through the prescribed spatial widths
\[
  r_0>r_1>r_2>\cdots\longrightarrow0
\]
at times for which \(c r_n\leq\rho(t_n)\leq C r_n\), with constants \(c,C>0\) independent of \(n\), and those times would have to obey
\[
  t_0<t_1<t_2<\cdots\uparrow T_*<\infty,
  \qquad
  t_{n+1}-t_n\asymp\tau_n,
\]
with the time-comparison constants also independent of \(n\). Only under these same-history requirements does the remaining time satisfy
\[
  T_*-t_n
  \asymp
  \sum_{k=n}^{\infty}\tau_k
  =
  \frac{4r_n^2}{3\nu}.
\]
At a time when the actual concentration width is comparable to \(r_n\), the next transition and the entire remaining interval are both of order \(r_n^2/\nu\). The same velocity field must produce each further narrowing inside an interval collapsing with the square of its current width.